\section{Introduction}

Financial literacy is a prerequisite for effective participation in modern financial markets. In India, the rapid expansion of digital financial infrastructure, mutual fund participation and online investment platforms has created both opportunities and risks for households. National surveys and the World Bank Global Findex 2025 report highlight persistent gaps in financial knowledge, particularly among first-time investors who face difficulties in understanding risk, fees, product structures and the long-term effects of compounding~\cite{worldbank2025}.

Beginner investors confront several interrelated challenges. First, financial terminology such as net asset value (NAV), tracking error, drawdown and riskometer ratings is often unfamiliar. Second, the sheer volume of available schemes, with over 38,000 registered mutual fund schemes tracked by AMFI, makes informed product selection infeasible without structured guidance. Third, risk awareness is frequently absent: investors may chase recent returns while underestimating volatility and downside exposure~\cite{sebi_understanding_mutual_funds}. Fourth, the absence of a safe experimentation environment means that costly mistakes are made with real capital.

These challenges are not merely individual; they bear directly on the Sustainable Development Goals. Quality education (SDG 4) requires accessible financial literacy resources. Decent work and economic growth (SDG 8) depend on informed personal financial decisions that support household stability. Reduced inequalities (SDG 10) are supported when financial education reduces barriers to understanding formal financial products. Responsible consumption and production (SDG 12) are encouraged when risk awareness and simulated experience promote considered financial decisions.

This paper presents WealthGenie, an integrated financial literacy and investment decision support platform that addresses these challenges through five coordinated modules. The platform is designed for Indian beginners and operates entirely in Indian Rupees with en-IN number formatting, India-specific instruments such as SIPs, PPF, EPF and DICGC-insured deposits, and content grounded in SEBI, RBI and AMFI publications.

The research contributions of this work are threefold. First, we present an integrated learning pipeline that moves users from structured financial literacy through personal assessment to explainable investment exploration, simulated trading and scenario comparison. Second, we describe a deterministic, risk-aware mutual fund recommendation engine that operates on verified AMFI and SEBI datasets and publishes its scoring weights and per-factor breakdowns, enabling full traceability of every recommendation. Third, we report implementation-consistency and explainability evidence for the deterministic assessment and Digital Twin components, evaluated against an independently implemented reference on synthetic profiles.

The remainder of this paper is organized as follows. Section II discusses SDG alignment. Section III describes the overall platform. Sections IV through VIII detail the five core modules. Section IX presents the system implementation and evaluation. Section X discusses limitations, and Section XI concludes the paper.